\documentclass[11pt]{article}

\usepackage[margin=1in]{geometry}
\usepackage{amsmath,amssymb}
\usepackage{booktabs}
\usepackage{array}
\usepackage{enumitem}
\usepackage[hidelinks]{hyperref}

\newcommand{\expit}{\operatorname{expit}}
\newcommand{\ind}{\mathbb{I}}
\newcommand{\E}{\mathbb{E}}
\newcommand{\Prb}{\operatorname{Pr}}
\newcommand{\RMST}{\operatorname{RMST}}

\title{ArraySim5-3: Site-Practice Heterogeneity\\
Data-Generating Process, Oracle Truth, and CCW Estimators}
\author{}
\date{}

\begin{document}
\maketitle

\section{Purpose and simulation structure}

ArraySim5-3 evaluates clone--censor--weight (CCW) estimation when treatment
initiation practices differ systematically across sites even after adjustment
for measured baseline characteristics and time-varying clinical history. The
simulation deliberately holds patient mix constant across sites. Therefore,
between-site differences in treatment initiation are generated by a residual
site-practice effect rather than by differences in the distribution of
measured patient characteristics.

There are three sites, indexed by $k\in\{1,2,3\}$, with 1,000 patients per
site. Patients are indexed by $i$, and discrete follow-up intervals are
indexed by $m=1,\ldots,M$, where $M=25$. The simulation grid is
\[
  \tau\in\{5,10,15\},\qquad
  \beta_{\mathrm{trt}}\in\{-1.0,-0.7,-0.5\},\qquad
  h\in\{\text{low},\text{moderate},\text{high}\},
\]
where $\tau$ is the grace-period length and $h$ indexes the magnitude of
site-practice heterogeneity. This gives $3\times3\times3=27$ simulation
cells. Each cell contains 100 Monte Carlo replicates.

\section{Notation}

For patient $i$ at site $k$, let
\begin{itemize}[leftmargin=2em]
  \item $X_{1ik}$ and $X_{2ik}$ denote baseline covariates;
  \item $L_{1ikm}$ and $L_{2ikm}$ denote time-varying covariates;
  \item $S_{ik}$ denote the interval of treatment initiation;
  \item $T_{ik}$ denote the terminal-event interval;
  \item $A_{ik}^{\tau}=\ind(S_{ik}\leq\tau)$ indicate initiation within the
        grace period; and
  \item $g\in\{1,0\}$ index the strategies ``initiate within $\tau$'' and
        ``do not initiate within $\tau$,'' respectively.
\end{itemize}

If a patient does not initiate or experience the event during observed
follow-up, the corresponding time is coded as $M+1$.

\section{Data-generating process}

\subsection{Baseline patient mix}

The baseline distributions are identical at all three sites:
\[
  X_{1ik}\sim N(0,1),
  \qquad
  X_{2ik}\sim\operatorname{Bernoulli}(0.40).
\]
Thus,
\[
  (X_{1ik},X_{2ik})\mid K_i=k
  \overset{d}{=}
  (X_{1ik'},X_{2ik'})\mid K_i=k'
\]
for every pair of sites $k$ and $k'$. Site-practice heterogeneity is
therefore not induced by different measured patient populations.

\subsection{Time-varying covariates}

For $\ell\in\{1,2\}$, the time-varying covariates follow
\[
  L_{\ell,ikm}
  =0.6L_{\ell,ik,m-1}
   +0.3X_{1ik}
   +0.2X_{2ik}
   +\varepsilon_{\ell,ikm},
  \qquad
  \varepsilon_{\ell,ikm}\sim N(0,1),
\]
with $L_{\ell,ik0}=0$. Innovations are generated separately for $L_1$ and
$L_2$. Treatment does not affect the subsequent covariate process. Covariates
are generated only while a patient remains alive; values after a terminal
event are left undefined.

\subsection{Treatment initiation}

Among patients who are alive and have not initiated before interval $m$, let
\[
  q_{ikm}
  =\Prb(S_{ik}=m\mid S_{ik}\geq m,\,T_{ik}\geq m,
                         X_{ik},L_{ikm},K_i=k).
\]
The treatment-initiation DGP is
\begin{equation}
\label{eq:init-dgp}
  \operatorname{logit}(q_{ikm})
  =\alpha_{kh}
   +0.8X_{1ik}
   -0.3X_{2ik}
   +0.5L_{1ikm}
   +0.4L_{2ikm}.
\end{equation}
The scenario-specific site intercepts are
\begin{table}[ht]
\centering
\caption{Site-specific treatment-initiation intercepts.}
\begin{tabular}{lrrrr}
\toprule
Practice heterogeneity & Site 1 & Site 2 & Site 3 & Range \\
\midrule
Low      & $-3.25$ & $-3.00$ & $-2.75$ & $0.50$ \\
Moderate & $-3.75$ & $-3.00$ & $-2.25$ & $1.50$ \\
High     & $-4.50$ & $-3.00$ & $-1.50$ & $3.00$ \\
\bottomrule
\end{tabular}
\end{table}

Every scenario is centered at $-3$. Consequently, the scenario axis changes
the degree of between-site practice heterogeneity without uniformly shifting
the treatment-initiation intercept across all sites. At zero covariate values,
the corresponding per-interval probabilities are approximately
\begin{center}
\begin{tabular}{lrrr}
\toprule
Practice heterogeneity & Site 1 & Site 2 & Site 3 \\
\midrule
Low      & $3.7\%$ & $4.7\%$ & $6.0\%$ \\
Moderate & $2.3\%$ & $4.7\%$ & $9.5\%$ \\
High     & $1.1\%$ & $4.7\%$ & $18.2\%$ \\
\bottomrule
\end{tabular}
\end{center}
Thus, two patients with the same measured history may have different
initiation probabilities solely because they receive care at different sites.

\subsection{Terminal-event process}

Conditional on being alive at the start of interval $m$, define the event
hazard
\[
  p^Y_{ikm}=\Prb(T_{ik}=m\mid T_{ik}\geq m,\mathcal H_{ikm}),
\]
where $\mathcal H_{ikm}$ contains the observed covariate and treatment history.
The event DGP is
\begin{align}
\label{eq:event-dgp}
  \operatorname{logit}(p^Y_{ikm})
  ={}&-2.5
      +0.5X_{1ik}+0.4X_{2ik}
      +0.4L_{1ikm}+0.3L_{2ikm} \notag\\
     &+0.2L_{1ik,m-1}+0.15L_{2ik,m-1}
      +\beta_{\mathrm{trt}}
       \ind(S_{ik}\leq\tau,\,m\geq S_{ik}).
\end{align}
Treatment affects the event hazard only if initiation occurs within the grace
period and the current interval is at or after initiation. Initiation after
$\tau$ does not activate the treatment effect. There is no direct site effect
in Equation~\eqref{eq:event-dgp}; ArraySim5-3 isolates site heterogeneity in
treatment initiation rather than site heterogeneity in outcome risk.

Within each interval, the simulation first generates the current $L$ values,
then treatment initiation, and finally the terminal event. A patient may
therefore initiate and experience the event in the same interval, with
treatment initiation occurring first.

\section{CCW strategies and artificial censoring}

Each patient is conceptually cloned into two strategies:
\begin{align*}
  g=1:&\quad \text{initiate treatment within the grace period},\\
  g=0:&\quad \text{do not initiate treatment within the grace period}.
\end{align*}
The code implements these clones through strategy-specific artificial
censoring times rather than by physically duplicating the data.

For $g=1$, a patient who fails to initiate by $\tau$ deviates during interval
$\tau$ and is censored at the end of interval $\tau-1$:
\[
  C^1_{ik}
  =\begin{cases}
      M+1, & S_{ik}\leq\tau,\\
      \tau-1, & S_{ik}>\tau.
    \end{cases}
\]
For $g=0$, a patient who initiates during interval $S_{ik}\leq\tau$ deviates
during that interval and is censored at the end of the preceding interval:
\[
  C^0_{ik}
  =\begin{cases}
      S_{ik}-1, & S_{ik}\leq\tau,\\
      M+1, & S_{ik}>\tau.
    \end{cases}
\]
This convention excludes an event in the initiation interval from the
$g=0$ clone but allows it to contribute to the $g=1$ clone.

\section{Oracle truth}

\subsection{Oracle population}

For every combination of $(\tau,\beta_{\mathrm{trt}},h)$, an oracle cohort of
$N_{\mathrm{oracle}}=3{,}000{,}000$ patients is simulated. The cohort is
allocated equally across the three sites, so the target population has site
weights
\[
  \pi_k=\frac{1}{3},\qquad k=1,2,3.
\]
The oracle uses the scenario-specific site intercepts and the same DGP as the
finite samples. Common random-number seeds are used across grid cells to
reduce unnecessary Monte Carlo variation.

\subsection{True denominator and common numerator}

The oracle denominator is the true treatment-initiation probability from
Equation~\eqref{eq:init-dgp}, rather than a fitted regression:
\[
  q_{ikm}^{D,\mathrm{true}}
  =\expit\!\left(
      \alpha_{kh}+0.8X_{1ik}-0.3X_{2ik}
      +0.5L_{1ikm}+0.4L_{2ikm}
    \right).
\]
The common numerator hazard in the oracle population is
\begin{equation}
\label{eq:oracle-num}
  q_m^{N,\mathrm{oracle}}
  =
  \frac{
    \sum_{k=1}^3\sum_i
      \ind(S_{ik}=m,\,T_{ik}\geq m)
  }{
    \sum_{k=1}^3\sum_i
      \ind(S_{ik}\geq m,\,T_{ik}\geq m)
  }.
\end{equation}
Because treatment initiation is generated before the event within an
interval, a patient with $T_{ik}=m$ is eligible for the initiation decision
at interval $m$.

\subsection{Oracle weights}

Write $q_m^N=q_m^{N,\mathrm{oracle}}$ and
$q_{ikm}^D=q_{ikm}^{D,\mathrm{true}}$. Before initiation, the incremental
waiting factor is
\[
  r^{\mathrm{wait}}_{ikm}
  =\frac{1-q_m^N}{1-q_{ikm}^D},
\]
and at the initiation interval the incremental initiation factor is
\[
  r^{\mathrm{init}}_{ikm}
  =\frac{q_m^N}{q_{ikm}^D}.
\]
For the $g=1$ clone, the stabilized weight accumulates waiting factors before
initiation and an initiation factor at $S_{ik}$:
\begin{equation}
\label{eq:w1}
  W^1_{ikm}
  =\prod_{j=1}^{\min(m,\tau)}
    \left\{
      \left(\frac{q_j^N}{q_{ikj}^D}\right)^{\ind(S_{ik}=j)}
      \left(\frac{1-q_j^N}{1-q_{ikj}^D}\right)^{\ind(S_{ik}>j)}
    \right\},
\end{equation}
where factors are evaluated only while the patient remains alive and
eligible for initiation. For the $g=0$ clone,
\begin{equation}
\label{eq:w0}
  W^0_{ikm}
  =\prod_{j=1}^{\min(m,\tau,S_{ik}-1)}
    \frac{1-q_j^N}{1-q_{ikj}^D}.
\end{equation}
After $\tau$, weights are carried forward without further updating. The
current simulation uses quantile limits $(0,1)$, which results in no effective
weight truncation.

\subsection{Oracle survival and target quantities}

For strategy $g$, define
\[
  R^g_{ikm}
  =\ind(T_{ik}>m-1,\,C^g_{ik}>m-1)
\]
and
\[
  D^g_{ikm}
  =\ind(R^g_{ikm}=1,\,T_{ik}=m,\,T_{ik}\leq C^g_{ik}).
\]
The oracle weighted hazard and survival curve are
\[
  \lambda^{g,\mathrm{oracle}}_m
  =\frac{\sum_{k,i}W^g_{ikm}D^g_{ikm}}
         {\sum_{k,i}W^g_{ikm}R^g_{ikm}},
  \qquad
  S^{g,\mathrm{oracle}}(m)
  =\prod_{j=1}^m\left(1-\lambda^{g,\mathrm{oracle}}_j\right).
\]
The oracle therefore defines the Monte Carlo truth for the CCW estimand under
the equal-site target population and pooled marginal initiation-timing
distribution represented by the common numerator.

At $M=25$, define
\[
  \psi_g=1-S^g(M),
  \qquad
  \RMST_g=\sum_{m=0}^{M-1}S^g(m).
\]
The reported contrasts are
\[
  \mathrm{RD}=\psi_1-\psi_0,
  \qquad
  \mathrm{RR}=\frac{\psi_1}{\psi_0},
\]
\[
  \mathrm{OR}
  =\frac{\psi_1/(1-\psi_1)}{\psi_0/(1-\psi_0)},
  \qquad
  \RMST_{\mathrm{diff}}=\RMST_1-\RMST_0.
\]

\section{Federated CCW}

\subsection{Site-specific denominator nuisance models}

Each site constructs a person-interval dataset containing intervals in which
the patient is alive and has not initiated before the interval. Site $k$ fits
the logistic model
\begin{align}
\label{eq:fitted-den}
  \operatorname{logit}(\widehat q^D_{ikm})
  ={}&\widehat\gamma_{km}
      +\widehat\beta_{1k}X_{1ik}
      +\widehat\beta_{2k}X_{2ik}
      +\widehat\beta_{3k}L_{1ikm}
      +\widehat\beta_{4k}L_{2ikm} \notag\\
     &+\widehat\beta_{5k}L_{1ik,m-1}
      +\widehat\beta_{6k}L_{2ik,m-1}.
\end{align}
The interval effect $\widehat\gamma_{km}$ is represented by a factor for
$m$. Equation~\eqref{eq:fitted-den} is more flexible than the initiation DGP:
it allows a separate baseline initiation hazard by interval and includes
lagged $L$ terms whose coefficients are zero in the true initiation model.

All coefficients in Equation~\eqref{eq:fitted-den} are estimated separately
within each site. There is no cross-site sharing or partial pooling of the
denominator regression coefficients.

\subsection{Federated common numerator}

Site $k$ releases the interval-level counts
\[
  N^{\mathrm{init}}_{km}
  =\sum_i\ind(S_{ik}=m,\,T_{ik}\geq m)
\]
and
\[
  N^{\mathrm{eligible}}_{km}
  =\sum_i\ind(S_{ik}\geq m,\,T_{ik}\geq m).
\]
The central server forms
\begin{equation}
\label{eq:fed-num}
  \widehat q_m^N
  =\frac{\sum_k N^{\mathrm{init}}_{km}}
         {\sum_k N^{\mathrm{eligible}}_{km}}
\end{equation}
and broadcasts this common numerator vector to all sites. Each site then
constructs Equations~\eqref{eq:w1} and \eqref{eq:w0}, replacing the oracle
quantities with $\widehat q_m^N$ and $\widehat q^D_{ikm}$.

\subsection{Local weighted sufficient statistics}

Using its local individual-level records, site $k$ computes
\[
  DW^g_{km}=\sum_i W^g_{ikm}D^g_{ikm},
  \qquad
  RW^g_{km}=\sum_i W^g_{ikm}R^g_{ikm}.
\]
Only these weighted event and risk totals are required centrally for point
estimation. The central server calculates
\begin{equation}
\label{eq:fed-hazard}
  \widehat\lambda^{g,\mathrm{fed}}_m
  =\frac{\sum_k DW^g_{km}}{\sum_k RW^g_{km}}
\end{equation}
and
\[
  \widehat S^{g,\mathrm{fed}}(m)
  =\prod_{j=1}^m
     \left(1-\widehat\lambda^{g,\mathrm{fed}}_j\right).
\]
The treatment-effect estimates are obtained by substituting these survival
curves into the RD, RR, OR, and RMST formulas above. This is a pooled weighted
survival estimator: site-specific effect estimates are not calculated and
then meta-analyzed.

\subsection{Federated variance calculation}

After computing the global hazards and survival curves, the central server
returns the quantities needed for a second local pass. A strategy-specific
incremental influence contribution is
\[
  U^g_{ikm}
  =\frac{1}{1-\widehat\lambda^g_m}
   \frac{W^g_{ikm}}{\sum_k RW^g_{km}}
   \left(D^g_{ikm}-\widehat\lambda^g_mR^g_{ikm}\right).
\]
Sites calculate individual influence-function contributions and release sums
of their squares. For example,
\[
  \widehat{\operatorname{SE}}(\widehat{\mathrm{RD}})
  =\left\{
      \sum_k\sum_i
      \left(IF^{\psi_1}_{ik}-IF^{\psi_0}_{ik}\right)^2
    \right\}^{1/2}.
\]
Analogous calculations are used for log-RR, log-OR, and RMST contrasts. The
implemented variance calculation treats the estimated weights as fixed and
does not add nuisance-model estimation uncertainty.

\section{Aligned fully site-stratified pooled CCW}

The pooled comparator has access to all individual-level records at a central
location, but it deliberately uses the same statistical specification as the
federated estimator. Specifically, it
\begin{enumerate}[leftmargin=2em]
  \item separates the pooled data by site;
  \item fits Equation~\eqref{eq:fitted-den} separately within every site;
  \item constructs the same common numerator in Equation~\eqref{eq:fed-num};
  \item applies the same artificial censoring and weight construction;
  \item pools the weighted outcome records; and
  \item applies the same survival and influence-function calculations.
\end{enumerate}
It is not a conventional pooled logistic model with a single common slope
vector and site fixed effects. It is a fully site-stratified individual-data
implementation of the same estimator used by Fed-CCW.

The pooled weighted hazard is
\begin{equation}
\label{eq:pool-hazard}
  \widehat\lambda^{g,\mathrm{pool}}_m
  =\frac{
      \sum_k\sum_i W^g_{ikm}D^g_{ikm}
    }{
      \sum_k\sum_i W^g_{ikm}R^g_{ikm}
    }.
\end{equation}
Because
\[
  DW^g_{km}=\sum_iW^g_{ikm}D^g_{ikm},
  \qquad
  RW^g_{km}=\sum_iW^g_{ikm}R^g_{ikm},
\]
Equations~\eqref{eq:fed-hazard} and \eqref{eq:pool-hazard} imply
\[
  \widehat\lambda^{g,\mathrm{fed}}_m
  =\widehat\lambda^{g,\mathrm{pool}}_m.
\]
It follows that
\[
  \widehat S^{g,\mathrm{fed}}(m)
  =\widehat S^{g,\mathrm{pool}}(m)
\]
and hence
\[
  \widehat\theta^{\mathrm{fed}}
  =\widehat\theta^{\mathrm{pool}}
\]
for every reported CCW estimand $\theta$. Because the pooled comparator also
uses the same influence-function decomposition, the implemented model-based
standard errors agree up to numerical precision.

The aligned pooled CCW is therefore an implementation gold standard rather
than a statistically different competing estimator. Its purpose is to show
that the federated exchange of common numerator counts, weighted event/risk
totals, and influence-function summaries reproduces the corresponding
individual-data analysis without transferring individual records.

\section{Where information is and is not shared}

The denominator nuisance regressions are fitted independently, so no site
borrows regression coefficients from another site. Cross-site information is
combined in two places:
\begin{enumerate}[leftmargin=2em]
  \item the common numerator $\widehat q_m^N$, estimated from aggregated
        initiation and eligibility counts; and
  \item the final outcome estimator, formed by aggregating weighted event,
        risk, and influence-function summaries.
\end{enumerate}
Consequently, ArraySim5-3 evaluates lossless federation for a fully
site-stratified nuisance-model specification. It does not evaluate a shared
slope, random-effect, or hierarchical nuisance model.

\section{Monte Carlo performance evaluation}

For each replicate $r$ and estimand $\theta$, performance is evaluated against
the cell-specific oracle truth:
\[
  \operatorname{Error}_r(\theta)=\widehat\theta_r-\theta_{\mathrm{oracle}},
\]
\[
  \operatorname{Bias}(\theta)
  =\frac{1}{R}\sum_{r=1}^R
    \left(\widehat\theta_r-\theta_{\mathrm{oracle}}\right),
\]
and
\[
  \operatorname{RMSE}(\theta)
  =\left\{
      \frac{1}{R}\sum_{r=1}^R
      \left(\widehat\theta_r-\theta_{\mathrm{oracle}}\right)^2
    \right\}^{1/2},
  \qquad R=100.
\]
Because Fed-CCW and the aligned pooled CCW implement the same estimator, their
replicate-level estimates, biases, RMSEs, and current model-based confidence
intervals are expected to overlap.

\end{document}
